
\documentclass[lettersize,journal]{IEEEtran}
\usepackage{amsmath,amsfonts}
% \usepackage{algorithmic}
\usepackage{array}
\usepackage[caption=false,font=normalsize,labelfont=sf,textfont=sf]{subfig}
\usepackage{textcomp}
\usepackage{stfloats}
\usepackage{url}
\usepackage{verbatim}
\usepackage{graphicx}
\usepackage{cite}
\hyphenation{op-tical net-works semi-conduc-tor IEEE-Xplore}

% Optional math commands from https://github.com/goodfeli/dlbook_notation.
%%%%% NEW MATH DEFINITIONS %%%%%

\usepackage{amsmath,amsfonts,bm}

% Mark sections of captions for referring to divisions of figures
\newcommand{\figleft}{{\em (Left)}}
\newcommand{\figcenter}{{\em (Center)}}
\newcommand{\figright}{{\em (Right)}}
\newcommand{\figtop}{{\em (Top)}}
\newcommand{\figbottom}{{\em (Bottom)}}
\newcommand{\captiona}{{\em (a)}}
\newcommand{\captionb}{{\em (b)}}
\newcommand{\captionc}{{\em (c)}}
\newcommand{\captiond}{{\em (d)}}

% Highlight a newly defined term
\newcommand{\newterm}[1]{{\bf #1}}


% Figure reference, lower-case.
\def\figref#1{figure~\ref{#1}}
% Figure reference, capital. For start of sentence
\def\Figref#1{Figure~\ref{#1}}
\def\twofigref#1#2{figures \ref{#1} and \ref{#2}}
\def\quadfigref#1#2#3#4{figures \ref{#1}, \ref{#2}, \ref{#3} and \ref{#4}}
% Section reference, lower-case.
\def\secref#1{section~\ref{#1}}
% Section reference, capital.
\def\Secref#1{Section~\ref{#1}}
% Reference to two sections.
\def\twosecrefs#1#2{sections \ref{#1} and \ref{#2}}
% Reference to three sections.
\def\secrefs#1#2#3{sections \ref{#1}, \ref{#2} and \ref{#3}}
% Reference to an equation, lower-case.
\def\eqref#1{equation~\ref{#1}}
% Reference to an equation, upper case
\def\Eqref#1{Equation~\ref{#1}}
% A raw reference to an equation---avoid using if possible
\def\plaineqref#1{\ref{#1}}
% Reference to a chapter, lower-case.
\def\chapref#1{chapter~\ref{#1}}
% Reference to an equation, upper case.
\def\Chapref#1{Chapter~\ref{#1}}
% Reference to a range of chapters
\def\rangechapref#1#2{chapters\ref{#1}--\ref{#2}}
% Reference to an algorithm, lower-case.
\def\algref#1{algorithm~\ref{#1}}
% Reference to an algorithm, upper case.
\def\Algref#1{Algorithm~\ref{#1}}
\def\twoalgref#1#2{algorithms \ref{#1} and \ref{#2}}
\def\Twoalgref#1#2{Algorithms \ref{#1} and \ref{#2}}
% Reference to a part, lower case
\def\partref#1{part~\ref{#1}}
% Reference to a part, upper case
\def\Partref#1{Part~\ref{#1}}
\def\twopartref#1#2{parts \ref{#1} and \ref{#2}}

\def\ceil#1{\lceil #1 \rceil}
\def\floor#1{\lfloor #1 \rfloor}
\def\1{\bm{1}}
\newcommand{\train}{\mathcal{D}}
\newcommand{\valid}{\mathcal{D_{\mathrm{valid}}}}
\newcommand{\test}{\mathcal{D_{\mathrm{test}}}}

\def\eps{{\epsilon}}


% Random variables
\def\reta{{\textnormal{$\eta$}}}
\def\ra{{\textnormal{a}}}
\def\rb{{\textnormal{b}}}
\def\rc{{\textnormal{c}}}
\def\rd{{\textnormal{d}}}
\def\re{{\textnormal{e}}}
\def\rf{{\textnormal{f}}}
\def\rg{{\textnormal{g}}}
\def\rh{{\textnormal{h}}}
\def\ri{{\textnormal{i}}}
\def\rj{{\textnormal{j}}}
\def\rk{{\textnormal{k}}}
\def\rl{{\textnormal{l}}}
% rm is already a command, just don't name any random variables m
\def\rn{{\textnormal{n}}}
\def\ro{{\textnormal{o}}}
\def\rp{{\textnormal{p}}}
\def\rq{{\textnormal{q}}}
\def\rr{{\textnormal{r}}}
\def\rs{{\textnormal{s}}}
\def\rt{{\textnormal{t}}}
\def\ru{{\textnormal{u}}}
\def\rv{{\textnormal{v}}}
\def\rw{{\textnormal{w}}}
\def\rx{{\textnormal{x}}}
\def\ry{{\textnormal{y}}}
\def\rz{{\textnormal{z}}}

% Random vectors
\def\rvepsilon{{\mathbf{\epsilon}}}
\def\rvtheta{{\mathbf{\theta}}}
\def\rva{{\mathbf{a}}}
\def\rvb{{\mathbf{b}}}
\def\rvc{{\mathbf{c}}}
\def\rvd{{\mathbf{d}}}
\def\rve{{\mathbf{e}}}
\def\rvf{{\mathbf{f}}}
\def\rvg{{\mathbf{g}}}
\def\rvh{{\mathbf{h}}}
\def\rvu{{\mathbf{i}}}
\def\rvj{{\mathbf{j}}}
\def\rvk{{\mathbf{k}}}
\def\rvl{{\mathbf{l}}}
\def\rvm{{\mathbf{m}}}
\def\rvn{{\mathbf{n}}}
\def\rvo{{\mathbf{o}}}
\def\rvp{{\mathbf{p}}}
\def\rvq{{\mathbf{q}}}
\def\rvr{{\mathbf{r}}}
\def\rvs{{\mathbf{s}}}
\def\rvt{{\mathbf{t}}}
\def\rvu{{\mathbf{u}}}
\def\rvv{{\mathbf{v}}}
\def\rvw{{\mathbf{w}}}
\def\rvx{{\mathbf{x}}}
\def\rvy{{\mathbf{y}}}
\def\rvz{{\mathbf{z}}}

% Elements of random vectors
\def\erva{{\textnormal{a}}}
\def\ervb{{\textnormal{b}}}
\def\ervc{{\textnormal{c}}}
\def\ervd{{\textnormal{d}}}
\def\erve{{\textnormal{e}}}
\def\ervf{{\textnormal{f}}}
\def\ervg{{\textnormal{g}}}
\def\ervh{{\textnormal{h}}}
\def\ervi{{\textnormal{i}}}
\def\ervj{{\textnormal{j}}}
\def\ervk{{\textnormal{k}}}
\def\ervl{{\textnormal{l}}}
\def\ervm{{\textnormal{m}}}
\def\ervn{{\textnormal{n}}}
\def\ervo{{\textnormal{o}}}
\def\ervp{{\textnormal{p}}}
\def\ervq{{\textnormal{q}}}
\def\ervr{{\textnormal{r}}}
\def\ervs{{\textnormal{s}}}
\def\ervt{{\textnormal{t}}}
\def\ervu{{\textnormal{u}}}
\def\ervv{{\textnormal{v}}}
\def\ervw{{\textnormal{w}}}
\def\ervx{{\textnormal{x}}}
\def\ervy{{\textnormal{y}}}
\def\ervz{{\textnormal{z}}}

% Random matrices
\def\rmA{{\mathbf{A}}}
\def\rmB{{\mathbf{B}}}
\def\rmC{{\mathbf{C}}}
\def\rmD{{\mathbf{D}}}
\def\rmE{{\mathbf{E}}}
\def\rmF{{\mathbf{F}}}
\def\rmG{{\mathbf{G}}}
\def\rmH{{\mathbf{H}}}
\def\rmI{{\mathbf{I}}}
\def\rmJ{{\mathbf{J}}}
\def\rmK{{\mathbf{K}}}
\def\rmL{{\mathbf{L}}}
\def\rmM{{\mathbf{M}}}
\def\rmN{{\mathbf{N}}}
\def\rmO{{\mathbf{O}}}
\def\rmP{{\mathbf{P}}}
\def\rmQ{{\mathbf{Q}}}
\def\rmR{{\mathbf{R}}}
\def\rmS{{\mathbf{S}}}
\def\rmT{{\mathbf{T}}}
\def\rmU{{\mathbf{U}}}
\def\rmV{{\mathbf{V}}}
\def\rmW{{\mathbf{W}}}
\def\rmX{{\mathbf{X}}}
\def\rmY{{\mathbf{Y}}}
\def\rmZ{{\mathbf{Z}}}

% Elements of random matrices
\def\ermA{{\textnormal{A}}}
\def\ermB{{\textnormal{B}}}
\def\ermC{{\textnormal{C}}}
\def\ermD{{\textnormal{D}}}
\def\ermE{{\textnormal{E}}}
\def\ermF{{\textnormal{F}}}
\def\ermG{{\textnormal{G}}}
\def\ermH{{\textnormal{H}}}
\def\ermI{{\textnormal{I}}}
\def\ermJ{{\textnormal{J}}}
\def\ermK{{\textnormal{K}}}
\def\ermL{{\textnormal{L}}}
\def\ermM{{\textnormal{M}}}
\def\ermN{{\textnormal{N}}}
\def\ermO{{\textnormal{O}}}
\def\ermP{{\textnormal{P}}}
\def\ermQ{{\textnormal{Q}}}
\def\ermR{{\textnormal{R}}}
\def\ermS{{\textnormal{S}}}
\def\ermT{{\textnormal{T}}}
\def\ermU{{\textnormal{U}}}
\def\ermV{{\textnormal{V}}}
\def\ermW{{\textnormal{W}}}
\def\ermX{{\textnormal{X}}}
\def\ermY{{\textnormal{Y}}}
\def\ermZ{{\textnormal{Z}}}

% Vectors
\def\vzero{{\bm{0}}}
\def\vone{{\bm{1}}}
\def\vmu{{\bm{\mu}}}
\def\vtheta{{\bm{\theta}}}
\def\va{{\bm{a}}}
\def\vb{{\bm{b}}}
\def\vc{{\bm{c}}}
\def\vd{{\bm{d}}}
\def\ve{{\bm{e}}}
\def\vf{{\bm{f}}}
\def\vg{{\bm{g}}}
\def\vh{{\bm{h}}}
\def\vi{{\bm{i}}}
\def\vj{{\bm{j}}}
\def\vk{{\bm{k}}}
\def\vl{{\bm{l}}}
\def\vm{{\bm{m}}}
\def\vn{{\bm{n}}}
\def\vo{{\bm{o}}}
\def\vp{{\bm{p}}}
\def\vq{{\bm{q}}}
\def\vr{{\bm{r}}}
\def\vs{{\bm{s}}}
\def\vt{{\bm{t}}}
\def\vu{{\bm{u}}}
\def\vv{{\bm{v}}}
\def\vw{{\bm{w}}}
\def\vx{{\bm{x}}}
\def\vy{{\bm{y}}}
\def\vz{{\bm{z}}}

% Elements of vectors
\def\evalpha{{\alpha}}
\def\evbeta{{\beta}}
\def\evepsilon{{\epsilon}}
\def\evlambda{{\lambda}}
\def\evomega{{\omega}}
\def\evmu{{\mu}}
\def\evpsi{{\psi}}
\def\evsigma{{\sigma}}
\def\evtheta{{\theta}}
\def\eva{{a}}
\def\evb{{b}}
\def\evc{{c}}
\def\evd{{d}}
\def\eve{{e}}
\def\evf{{f}}
\def\evg{{g}}
\def\evh{{h}}
\def\evi{{i}}
\def\evj{{j}}
\def\evk{{k}}
\def\evl{{l}}
\def\evm{{m}}
\def\evn{{n}}
\def\evo{{o}}
\def\evp{{p}}
\def\evq{{q}}
\def\evr{{r}}
\def\evs{{s}}
\def\evt{{t}}
\def\evu{{u}}
\def\evv{{v}}
\def\evw{{w}}
\def\evx{{x}}
\def\evy{{y}}
\def\evz{{z}}

% Matrix
\def\mA{{\bm{A}}}
\def\mB{{\bm{B}}}
\def\mC{{\bm{C}}}
\def\mD{{\bm{D}}}
\def\mE{{\bm{E}}}
\def\mF{{\bm{F}}}
\def\mG{{\bm{G}}}
\def\mH{{\bm{H}}}
\def\mI{{\bm{I}}}
\def\mJ{{\bm{J}}}
\def\mK{{\bm{K}}}
\def\mL{{\bm{L}}}
\def\mM{{\bm{M}}}
\def\mN{{\bm{N}}}
\def\mO{{\bm{O}}}
\def\mP{{\bm{P}}}
\def\mQ{{\bm{Q}}}
\def\mR{{\bm{R}}}
\def\mS{{\bm{S}}}
\def\mT{{\bm{T}}}
\def\mU{{\bm{U}}}
\def\mV{{\bm{V}}}
\def\mW{{\bm{W}}}
\def\mX{{\bm{X}}}
\def\mY{{\bm{Y}}}
\def\mZ{{\bm{Z}}}
\def\mBeta{{\bm{\beta}}}
\def\mPhi{{\bm{\Phi}}}
\def\mLambda{{\bm{\Lambda}}}
\def\mSigma{{\bm{\Sigma}}}

% Tensor
\DeclareMathAlphabet{\mathsfit}{\encodingdefault}{\sfdefault}{m}{sl}
\SetMathAlphabet{\mathsfit}{bold}{\encodingdefault}{\sfdefault}{bx}{n}
\newcommand{\tens}[1]{\bm{\mathsfit{#1}}}
\def\tA{{\tens{A}}}
\def\tB{{\tens{B}}}
\def\tC{{\tens{C}}}
\def\tD{{\tens{D}}}
\def\tE{{\tens{E}}}
\def\tF{{\tens{F}}}
\def\tG{{\tens{G}}}
\def\tH{{\tens{H}}}
\def\tI{{\tens{I}}}
\def\tJ{{\tens{J}}}
\def\tK{{\tens{K}}}
\def\tL{{\tens{L}}}
\def\tM{{\tens{M}}}
\def\tN{{\tens{N}}}
\def\tO{{\tens{O}}}
\def\tP{{\tens{P}}}
\def\tQ{{\tens{Q}}}
\def\tR{{\tens{R}}}
\def\tS{{\tens{S}}}
\def\tT{{\tens{T}}}
\def\tU{{\tens{U}}}
\def\tV{{\tens{V}}}
\def\tW{{\tens{W}}}
\def\tX{{\tens{X}}}
\def\tY{{\tens{Y}}}
\def\tZ{{\tens{Z}}}


% Graph
\def\gA{{\mathcal{A}}}
\def\gB{{\mathcal{B}}}
\def\gC{{\mathcal{C}}}
\def\gD{{\mathcal{D}}}
\def\gE{{\mathcal{E}}}
\def\gF{{\mathcal{F}}}
\def\gG{{\mathcal{G}}}
\def\gH{{\mathcal{H}}}
\def\gI{{\mathcal{I}}}
\def\gJ{{\mathcal{J}}}
\def\gK{{\mathcal{K}}}
\def\gL{{\mathcal{L}}}
\def\gM{{\mathcal{M}}}
\def\gN{{\mathcal{N}}}
\def\gO{{\mathcal{O}}}
\def\gP{{\mathcal{P}}}
\def\gQ{{\mathcal{Q}}}
\def\gR{{\mathcal{R}}}
\def\gS{{\mathcal{S}}}
\def\gT{{\mathcal{T}}}
\def\gU{{\mathcal{U}}}
\def\gV{{\mathcal{V}}}
\def\gW{{\mathcal{W}}}
\def\gX{{\mathcal{X}}}
\def\gY{{\mathcal{Y}}}
\def\gZ{{\mathcal{Z}}}

% Sets
\def\sA{{\mathbb{A}}}
\def\sB{{\mathbb{B}}}
\def\sC{{\mathbb{C}}}
\def\sD{{\mathbb{D}}}
% Don't use a set called E, because this would be the same as our symbol
% for expectation.
\def\sF{{\mathbb{F}}}
\def\sG{{\mathbb{G}}}
\def\sH{{\mathbb{H}}}
\def\sI{{\mathbb{I}}}
\def\sJ{{\mathbb{J}}}
\def\sK{{\mathbb{K}}}
\def\sL{{\mathbb{L}}}
\def\sM{{\mathbb{M}}}
\def\sN{{\mathbb{N}}}
\def\sO{{\mathbb{O}}}
\def\sP{{\mathbb{P}}}
\def\sQ{{\mathbb{Q}}}
\def\sR{{\mathbb{R}}}
\def\sS{{\mathbb{S}}}
\def\sT{{\mathbb{T}}}
\def\sU{{\mathbb{U}}}
\def\sV{{\mathbb{V}}}
\def\sW{{\mathbb{W}}}
\def\sX{{\mathbb{X}}}
\def\sY{{\mathbb{Y}}}
\def\sZ{{\mathbb{Z}}}

% Entries of a matrix
\def\emLambda{{\Lambda}}
\def\emA{{A}}
\def\emB{{B}}
\def\emC{{C}}
\def\emD{{D}}
\def\emE{{E}}
\def\emF{{F}}
\def\emG{{G}}
\def\emH{{H}}
\def\emI{{I}}
\def\emJ{{J}}
\def\emK{{K}}
\def\emL{{L}}
\def\emM{{M}}
\def\emN{{N}}
\def\emO{{O}}
\def\emP{{P}}
\def\emQ{{Q}}
\def\emR{{R}}
\def\emS{{S}}
\def\emT{{T}}
\def\emU{{U}}
\def\emV{{V}}
\def\emW{{W}}
\def\emX{{X}}
\def\emY{{Y}}
\def\emZ{{Z}}
\def\emSigma{{\Sigma}}

% entries of a tensor
% Same font as tensor, without \bm wrapper
\newcommand{\etens}[1]{\mathsfit{#1}}
\def\etLambda{{\etens{\Lambda}}}
\def\etA{{\etens{A}}}
\def\etB{{\etens{B}}}
\def\etC{{\etens{C}}}
\def\etD{{\etens{D}}}
\def\etE{{\etens{E}}}
\def\etF{{\etens{F}}}
\def\etG{{\etens{G}}}
\def\etH{{\etens{H}}}
\def\etI{{\etens{I}}}
\def\etJ{{\etens{J}}}
\def\etK{{\etens{K}}}
\def\etL{{\etens{L}}}
\def\etM{{\etens{M}}}
\def\etN{{\etens{N}}}
\def\etO{{\etens{O}}}
\def\etP{{\etens{P}}}
\def\etQ{{\etens{Q}}}
\def\etR{{\etens{R}}}
\def\etS{{\etens{S}}}
\def\etT{{\etens{T}}}
\def\etU{{\etens{U}}}
\def\etV{{\etens{V}}}
\def\etW{{\etens{W}}}
\def\etX{{\etens{X}}}
\def\etY{{\etens{Y}}}
\def\etZ{{\etens{Z}}}

% The true underlying data generating distribution
\newcommand{\pdata}{p_{\rm{data}}}
% The empirical distribution defined by the training set
\newcommand{\ptrain}{\hat{p}_{\rm{data}}}
\newcommand{\Ptrain}{\hat{P}_{\rm{data}}}
% The model distribution
\newcommand{\pmodel}{p_{\rm{model}}}
\newcommand{\Pmodel}{P_{\rm{model}}}
\newcommand{\ptildemodel}{\tilde{p}_{\rm{model}}}
% Stochastic autoencoder distributions
\newcommand{\pencode}{p_{\rm{encoder}}}
\newcommand{\pdecode}{p_{\rm{decoder}}}
\newcommand{\precons}{p_{\rm{reconstruct}}}

\newcommand{\laplace}{\mathrm{Laplace}} % Laplace distribution

\newcommand{\E}{\mathbb{E}}
\newcommand{\Ls}{\mathcal{L}}
\newcommand{\R}{\mathbb{R}}
\newcommand{\emp}{\tilde{p}}
\newcommand{\lr}{\alpha}
\newcommand{\reg}{\lambda}
\newcommand{\rect}{\mathrm{rectifier}}
\newcommand{\softmax}{\mathrm{softmax}}
\newcommand{\sigmoid}{\sigma}
\newcommand{\softplus}{\zeta}
\newcommand{\KL}{D_{\mathrm{KL}}}
\newcommand{\Var}{\mathrm{Var}}
\newcommand{\standarderror}{\mathrm{SE}}
\newcommand{\Cov}{\mathrm{Cov}}
% Wolfram Mathworld says $L^2$ is for function spaces and $\ell^2$ is for vectors
% But then they seem to use $L^2$ for vectors throughout the site, and so does
% wikipedia.
\newcommand{\normlzero}{L^0}
\newcommand{\normlone}{L^1}
\newcommand{\normltwo}{L^2}
\newcommand{\normlp}{L^p}
\newcommand{\normmax}{L^\infty}

\newcommand{\parents}{Pa} % See usage in notation.tex. Chosen to match Daphne's book.

\DeclareMathOperator*{\argmax}{arg\,max}
\DeclareMathOperator*{\argmin}{arg\,min}

\DeclareMathOperator{\sign}{sign}
\DeclareMathOperator{\Tr}{Tr}
\let\ab\allowbreak


\usepackage{hyperref}
\usepackage{wrapfig}  
\graphicspath{ {./images/} }

% Required packages
\usepackage{algorithm}
\usepackage{algpseudocode}
\usepackage{amsmath} % for math environments
\usepackage{xcolor, colortbl}  % for color
\usepackage{tcolorbox} % for colored boxes
% Define custom colors
\definecolor{myblue}{RGB}{44, 130, 201}
\definecolor{mygreen}{RGB}{34, 177, 76}
\definecolor{myred}{RGB}{237, 28, 36}
\definecolor{myorange}{RGB}{255, 127, 39}

%for table
\usepackage{tabularx}
\usepackage{booktabs}
\usepackage{pifont}
\usepackage{adjustbox}
\usepackage{multirow}

%for bar plot
\usepackage{pgfplots}
\usepackage{tikz}
\pgfplotsset{compat=1.17}

% Define custom commands
\newcommand{\todo}[1]{\textcolor{red}{[#1]}}
\newcommand{\todoa}[1]{\textcolor{blue}{[#1]}}
\newcommand{\tohl}[1]{\textcolor{cyan}{[#1]}}

\newtheorem{assumption}{Assumption}
\newtheorem{theorem}{Theorem}
\newtheorem{lemma}{Lemma}
\newtheorem{proof}{Proof}
\newtheorem{definition}{Definition}

% Authors must not appear in the submitted version. They should be hidden
% as long as the \iclrfinalcopy macro remains commented out below.
% Non-anonymous submissions will be rejected without review.

\title{\textsf{ProtoGen}: Federated Privacy-Preserving Prototype Learning via Feature Generation in Heterogeneous Environments\\ (Supplementary Resource)}
\author{Sujit~Chowdhury, Aritra Bhaduri, Raju Halder% <-this % stops a space
\IEEEcompsocitemizethanks{\IEEEcompsocthanksitem Sujit Chowdhury, Aritra Bhaduri, Raju Halder are with the Department
of Computer Science and Engineering, Indian Institute of Technology Patna, 
India, 801106.\protect\\
% note need leading \protect in front of \\ to get a newline within \thanks as
% \\ is fragile and will error, could use \hfil\break instead.
E-mail: \{sujit\_2021cs35, aritra\_2101ai40, halder\}@iitp.ac.in}
% \And
% Ji Q. Ren \& Yevgeny LeNet \\
% Department of Computational Neuroscience \\
% University of the Witwatersrand \\
% Joburg, South Africa \\
% \texttt{\{robot,net\}@wits.ac.za} \\
% \AND
% Coauthor \\
% Affiliation \\
% Address \\
% \texttt{email}
}

% The \author macro works with any number of authors. There are two commands
% used to separate the names and addresses of multiple authors: \And and \AND.
%
% Using \And between authors leaves it to \LaTeX{} to determine where to break
% the lines. Using \AND forces a linebreak at that point. So, if \LaTeX{}
% puts 3 of 4 authors names on the first line, and the last on the second
% line, try using \AND instead of \And before the third author name.

\newcommand{\fix}{\marginpar{FIX}}
\newcommand{\new}{\marginpar{NEW}}

% \iclrfinalcopy % Uncomment for camera-ready version, but NOT for submission.
\begin{document}


\maketitle

\section{Theoretical Analysis of \textsf{ProtoGen}}\label{sec-th}

This section presents the following theorems, which collectively demonstrate the theoretical foundations of \textsf{ProtoGen}.

\subsection{Preliminaries and Notation}

Let $\sK$ be the set of all clients, with $|\sK| = K$. For each client $k \in \sK$, let $\sD_k$ be its local dataset, with $n_k = |\sD_k|$ and $n = \sum_{k=1}^K n_k$. Let $\theta_{k,t} \in \mathbb{R}^d$ denote the local model parameters of client $k$ at round $t$. Let $\rho_{k,t} \in \mathbb{R}^p$ and $\varrho_{k,t} \in \mathbb{R}^q$ denote the \emph{PPGM} encoder and decoder parameters of client $k$ at round $t$, respectively. $\psi_k \leftarrow\{\rho_k,\varrho_k\}$ represents the general \emph{PPGM} model parameters of client $k$. 

We define the following:
\begin{itemize}
    \item $\mathcal{L}_C(\theta; x, y)$: loss function for client models
    % \item $\mathcal{L}_{\emph{GEN}}(\rho, \varrho; x, y)$: loss function for \textsf{PPGM}
    \item $\mathcal{L}_{\emph{GEN}}(\psi; x, y)$: loss function for \emph{PPGM}
    \item $\mathcal{T}(\theta_k) = \mathbb{E}_{(x,y)\sim\sD_k}[\mathcal{L}_C(\theta_k; x, y)]$: expected combined local objective loss for client $k$
    % \item $\mathcal{V}(\rho_k, \varrho_k) = \frac{1}{K}\sum_{k=1}^K \mathbb{E}_{(x,y)\sim\sD_k}[\mathcal{L}_{\emph{PPGM}}(\rho, \varrho; x, y)]$: expected feature generation loss for client $k$.
    \item $\mathcal{V}(\psi_k) = \frac{1}{K}\sum_{k=1}^K \mathbb{E}_{(x,y)\sim\sD_k}[\mathcal{L}_{\emph{GEN}}(\psi_k; x, y)]$: expected feature generation loss for client $k$.
\end{itemize}

\subsection{Assumptions}

We make the following assumptions:

\begin{assumption}[Smoothness]
\label{asm:1}
$\mathcal{T}(\theta_k)$ is $L$-smooth for all $k$: $\forall \theta_k, \theta'_k \in \mathbb{R}^d$,
\begin{equation}
    \|\nabla \mathcal{T}(\theta_k) - \nabla \mathcal{T}(\theta'_k)\| \leq L_1\|\theta_k - \theta'_k\|
\end{equation}
where $L_1 > 0$ is the Lipschitz constant for the local model.
\end{assumption}


\begin{assumption}[Strong Convexity]
\label{asm:2}
$\mathcal{T}(\theta_k)$ is $\mu$-strongly convex for all $k$: $\forall \theta_k, \theta'_k \in \mathbb{R}^d$,
\begin{equation}
    \mathcal{T}(\theta'_k) \geq \mathcal{T}(\theta_k) + \langle\nabla \mathcal{T}(\theta_k), \theta'_k - \theta_k\rangle + \frac{\mu}{2}\|\theta'_k - \theta_k\|^2
\end{equation}
\end{assumption}

\begin{assumption}[Bounded Variance]
\label{asm:3}
The variance of stochastic gradients is bounded for both client models and \emph{PPGM}, $\forall k \in \sK, \forall \theta_k \in \mathbb{R}^d$:
\begin{equation}
    \mathbb{E}_{(x,y)\sim\sD_k}[\|\nabla\mathcal{L}_C(\theta_k; x, y) - \nabla \mathcal{T}(\theta_k)\|^2] \leq \sigma_1^2
\end{equation}
% Similarly, $, \quad \forall k \in \sK, \forall \phi_k, \psi_k$
\begin{equation}
    \mathbb{E}_{(x,y)\sim \sD_k}[\|\nabla\mathcal{L}_{\emph{GEN}}(\psi_k; x, y) - \nabla \mathcal{V}(\psi_k)\|^2] \leq \sigma_2^2
\end{equation}
\end{assumption}

\begin{assumption}[\emph{PPGM} Lipschitz Continuity]
\label{asm:4}
The \emph{PPGM} model is Lipschitz continuous with respect to its parameters:
\begin{equation}
% \begin{split}    
    \|\nabla \mathcal{V}(\psi_k) - \nabla \mathcal{V}(\psi_k') \| \leq L_2(\|\psi_k - \psi_k'\|)
% \end{split}
\end{equation}
where $L_2 > 0$ is the Lipschitz constant for the \emph{PPGM} model.
\end{assumption}

\begin{assumption}[Gradient Dissimilarity]
\label{asm:5}
For each client \( k \in \sK \), assume that the loss function \( \mathcal{T}(\theta_k) \) is differentiable and bounded below, with \( \theta_{k,t} \) denoting the model parameters at round \( t \), and \( \theta_k^* \) being a minimizer of \( \mathcal{T}(\theta_k) \). Then, there exist constants \( G \geq 0 \) and \( \beta \geq 0 \) such that:
\begin{equation}
\frac{1}{K}\sum_{k=1}^K \|\nabla \mathcal{T}(\theta_{k,0}) - \nabla \mathcal{T}(\theta_k^*)\|^2 
\leq G^2 + 2\beta \left( \mathcal{T}(\theta_{k,0}) - \mathcal{T}(\theta_k^*) \right)
\end{equation}
where \( G^2 \) is a non-negative constant that upper bounds the inherent variance in the client gradients (even near local optima), \( \beta \) is a non-negative constant that quantifies how gradient differences grow with suboptimality (i.e., how far \( \theta_{k,t} \) is from \( \theta_k^* \) in terms of objective value).

\end{assumption}


\subsection{Convergence Analysis}
\begin{lemma}[One-step Progress for Client Models]
\label{lemma:1}
For client $k \in \sK$, gradient computed with learning rate $\eta_l$ over $E$ local epochs at round $t \geq 0$, its one-step convergence bound follows:
\begin{equation}
\begin{split}
    &\mathbb{E}[\|\theta_{k,t+1} - \theta_k^*\|^2] \leq (1-\mu\eta_l)\mathbb{E}[\|\theta_{k,t}  - \theta_k^*\|^2]\\
    & \qquad - 2\eta_l\mathbb{E}[\mathcal{T}(\theta_{k,t}) - \mathcal{T}(\theta_k^*)]\\ 
    & \qquad \quad + \eta_l^2\left(\sigma_1^2 + 6L\gamma_1 E\eta_l + 8(E-1)^2\gamma_1^2\right)
\end{split}
\end{equation}
where $\gamma_1 = \max\{\frac{8L_1}{\mu}, E\}$
\end{lemma}

\begin{proof}
Let $g_{k,t}$ be the average stochastic gradient computed by client $k$ at round $t$ over $E$ local epochs. The update rule for client $k$'s model parameters is $\theta_{k,t+1} = \theta_{k,t} - \eta_l g_{k,t}$. We begin by expanding $\|\theta_{k,t+1} - \theta_k^*\|^2$:

\begin{align*}
\|\theta_{k,t+1} - \theta_k^*\|^2 &= \|\theta_{k,t} - \eta_l g_{k,t} - \theta_k^*\|^2 \\
&= \|\theta_{k,t} - \theta_k^*\|^2 - 2\eta_l \langle g_{k,t}, \theta_{k,t} - \theta_k^* \rangle\\ 
& \qquad+ \eta_l^2 \|g_{k,t}\|^2
\end{align*}

Taking the expectation of both sides:

\begin{align*}
\mathbb{E}[\|\theta_{k,t+1} - \theta_k^*\|^2] &= \mathbb{E}[\|\theta_{k,t} - \theta_k^*\|^2] - 2\eta_l \mathbb{E}[\langle g_{k,t}, \theta_{k,t} - \theta_k^* \rangle]\\ 
& \qquad \qquad + \eta_l^2 \mathbb{E}[\|g_{k,t}\|^2]
\end{align*}

Using the strong convexity and smoothness of $T(\theta_k)$, we can bound the inner product term:

\begin{align*}
\langle g_{k,t}, \theta_{k,t} - \theta_k^* \rangle &\geq \mathcal{T}(\theta_{k,t}) - \mathcal{T}(\theta_k^*) + \frac{\mu}{2}\|\theta_{k,t} - \theta_k^*\|^2\\ 
&\qquad - \langle g_{k,t} - \nabla \mathcal{T}(\theta_{k,t}), \theta_{k,t} - \theta_k^* \rangle
\end{align*}

For the gradient norm, we use the smoothness property and the definition of $\gamma_1$:

\begin{equation}
\mathbb{E}[\|g_{k,t}\|^2] \leq 4L_1(\mathcal{T}(\theta_{k,t}) - \mathcal{T}(\theta_k^*)) + 2\sigma_1^2 + 2\gamma_1^2
\end{equation}

Combining these bounds and applying Cauchy-Schwarz inequality to the error term, we get:
\begin{align*}
\mathbb{E}[\|\theta_{k,t+1} - \theta_k^*\|^2] &\leq (1 - 2\mu\eta_l)\mathbb{E}[\|\theta_{k,t} - \theta_k^*\|^2]\\ 
&- 2\eta_l(1 - \eta_l L_1)\mathbb{E}[\mathcal{T}(\theta_{k,t}) - \mathcal{T}(\theta_k^*)] \\
&\quad + \eta_l^2(2\sigma_1^2 + 2\gamma_1^2) + \eta_l\sigma_1^2
\end{align*}

To account for $E$ local epochs, we model this as $E$ consecutive updates with learning rate $\eta_l$. Applying the above inequality $E$ times and using the convexity of $\mathcal{T}(\theta_k)$, we get the following:

\begin{align*}
% \centering
&\mathbb{E}[\|\theta_{k,t}^E - \theta_k^*\|^2] \leq (1 - 2\mu\eta_l)^E\mathbb{E}[\|\theta_{k,t} - \theta_k^*\|^2]\\ 
& \qquad \qquad - 2E\eta_l(1 - \eta_l L_1)\mathbb{E}[\mathcal{T}(\theta_{k,t}) - \mathcal{T}(\theta_k^*)] \\
&\qquad \qquad \quad + E\eta_l^2(2\sigma_1^2 + 2\gamma_1^2) + E\eta_l\sigma_1^2\\ 
&\qquad \qquad \qquad + 4L_1\eta_l^2E(E-1)\gamma_1^2
\end{align*}


Finally, using the fact that $\eta_l \leq 1/(4L_1)$ (which follows from our choice of $\eta_l$) and combining like terms, we obtain the stated bound:

\begin{equation}
\begin{split}
\mathbb{E}[\|\theta_{k,t+1} - \theta_k^*\|^2] \leq& (1-\mu\eta_l)\mathbb{E}[\|\theta_{k,t} - \theta_k^*\|^2]\\ 
&- 2\eta_l\mathbb{E}[\mathcal{T}(\theta_{k,t}) - \mathcal{T}(\theta_k^*)] \\
&+ \eta_l^2(\sigma_1^2 + 6L\gamma E\eta_l + 8(E-1)^2\gamma_1^2)
\end{split}
\end{equation}

\end{proof}

\begin{lemma}[One-step Progress for \emph{PPGM}]
\label{lemma:2}
For any round $t \geq 0$, the convergence bound for \emph{PPGM} follows:
\begin{equation}
\begin{split}
    \mathbb{E}[\|\psi_{k,t+1} - \psi_k^*\|^2]
    &\leq \mathbb{E}[\|\psi_{k,t} - \psi_k^*\|^2]\\
    &\quad - 2\eta_g \mathbb{E}[\mathcal{V}(\psi_{k,t}) - \mathcal{V}(\psi_k^*)] \\
    &\quad - \frac{\eta_g}{L_2} \sigma_2^2 + \eta_g^2 \sigma_2^2
\end{split}
\end{equation}
\end{lemma}


\begin{proof}
Let $g'_t = \nabla \mathcal{V}(\psi_{k,t})$ be the gradient of the \emph{PPGM} loss function at round $t$. The update rule for $\psi_k$ with learning rate $\eta_g$ defined as:
\begin{equation}
\psi_{k,t+1} = \psi_{k,t} - \eta_g g'_t
\end{equation}
% For simplicity, we denote $\mathbb{E}[\|\rho_{k,t+1} - \rho_k^*\|^2 + \|\varrho_{k,t+1} - \varrho_k^*\|^2]$ as $\|(\rho_{k,t+1}, \varrho_{k,t+1}) - (\rho_k^*, \varrho_k^*)\|^2$.\\
We begin by expanding $\|\psi_{k,t+1} - \psi_k^*\|^2$:

\begin{align*}
&\|\psi_{k,t+1} - \psi_k^*\|^2 = \|\psi_{k,t} - \eta_g g'_t - \psi_k^*\|^2 \\
&= \|\psi_{k,t} - \psi_k^*\|^2 - 2\eta_g \langle g'_t, \psi_{k,t} - \psi_k^* \rangle + \eta_g^2 \|g'_t\|^2
\end{align*}

Taking the expectation of both sides:

\begin{equation}
\label{eq:21}
\begin{split}
\mathbb{E}[\|\psi_{k,t+1} - \psi_k^*\|^2] &= \mathbb{E}[\|\psi_{k,t} - \psi_k^*\|^2] \\
&\qquad- 2\eta_g \mathbb{E}[\langle g'_t, \psi_{k,t} - \psi_k^* \rangle]\\ 
&\qquad \quad+ \eta_g^2 \mathbb{E}[\|g'_t\|^2]
\end{split}
\end{equation}

By the $L$-smoothness of $Y(\rho_k,\varrho_k)$ with Lipschitz constant $L_2 > 0$, we have:

\begin{equation}
\begin{split}
&\mathcal{V}(\psi_k^*) \geq \mathcal{V}(\psi_{k,t}) + \langle g'_t, \psi_k^* - \psi_{k,t} \rangle + \frac{1}{2L_2}\|g'_t\|^2    
\end{split}
\end{equation}

Rearranging this inequality:

\begin{equation}
\label{eq:23}
\begin{split}
&\langle g'_t, \psi_k^* - \psi_{k,t} \rangle \geq \mathcal{V}(\psi_{k,t})- \mathcal{V}(\psi_k^*) + \frac{1}{2L_2}\|g'_t\|^2
\end{split}
\end{equation}

Substituting Equation~\ref{eq:23} into Equation~\ref{eq:21}, we have:

\begin{align*}
&\mathbb{E}[\|\psi_{k,t+1} - \psi_k^*\|^2] \leq \mathbb{E}[\|\psi_{k,t} - \psi_k^*\|^2]\\
&\qquad\quad- 2\eta_g \mathbb{E}[\mathcal{V}(\psi_{k,t})- \mathcal{V}(\psi_k^*)]\\ 
&\qquad \qquad- \frac{\eta_g }{L_2}\mathbb{E}[\|g'_t\|^2]+ \eta_g^2 \mathbb{E}[\|g'_t\|^2]
\end{align*}

Now, using the assumption of bounded variance of stochastic gradients, $\mathbb{E}[\|g_t\|^2] \leq \sigma_2^2$, we get:

\begin{align*}
&\mathbb{E}[\|\psi_{k,t+1} - \psi_k^*\|^2] \leq \mathbb{E}[\|\psi_{k,t} - \psi_k^*\|^2]\\
&\qquad\quad- 2\eta_g \mathbb{E}[\mathcal{V}(\psi_{k,t})- \mathcal{V}(\psi_k^*)]\\ 
&\qquad\qquad - \frac{\eta_g}{L_2} \sigma_2^2 + \eta_g^2 \sigma_2^2
\end{align*}

\end{proof}

\begin{theorem}[Convergence of \textsf{ProtoGen} in Convex Settings]
\label{theo:1}
Let Assumptions 1-4 hold. Let the learning rates be set as $\eta_l = \frac{2}{\mu(t+\gamma_1)}$ for client models and $\eta_g = \frac{2}{L_2(t+\gamma_2)}$ for the \emph{PPGM}, where $\gamma_1 = \max\{\frac{8L_1}{\mu}, E\}$, $\gamma_2 = \max\{\frac{8L_2}{\mu}, E\}$, and $E$ is the number of local epochs. Then, for $T \geq 1$ number of global rounds of training, the output of \textsf{ProtoGen} satisfies, following convergence bound:
\begin{equation}
\begin{split}
    &\mathbb{E}[\mathcal{T}(\bar{\theta}_{k,T})] - \mathcal{T}(\theta_k^*)\\ 
    &\qquad \leq \frac{4L_1\gamma_1}{\mu T}\left(1 + \log\left(\frac{T}{\gamma_1} + 1\right)\right)
\end{split}
\end{equation}
\begin{equation}
\begin{split}
    &\mathbb{E}[\mathcal{V}(\bar{\psi}_{k,T})] - \mathcal{V}(\psi_k^*)\\ 
    & \qquad \leq \frac{4L_2\gamma_2}{T}\left(1 + \log\left(\frac{T}{\gamma_2} + 1\right)\right)
\end{split}
\end{equation}
where $\bar{\theta}_{k,T} = \frac{1}{T}\sum_{t=1}^T \theta_{k,t}$, $\bar{\psi}_T = \frac{1}{T}\sum_{t=1}^T \psi_t$, and $\theta_k^*, \psi^*$ are the respective optimal parameters.
\end{theorem}


\begin{proof}[Theorem~\ref{theo:1} (Convex Client Model Bound)]
We first derive the convergence bound for the local model of any client $k$. By Lemma~\ref{lemma:1}, summing over $T$ rounds gives:
\begin{equation}
\begin{split}
    \sum_{t=0}^{T-1} 2\eta_{l,t} \mathbb{E}[\mathcal{T}(\theta_{k,t}) - \mathcal{T}(\theta_k^*)]
    \leq \|\theta_{k,0} - \theta_k^*\|^2 + \sum_{t=0}^{T-1} \eta_{l,t}^2 C,
\end{split}
\end{equation}
where $C = \sigma_1^2 + 6L_1\gamma_1 E\eta_{l,t} + 8(E-1)^2\gamma_1^2$ and the learning rate is $\eta_{l,t} = \frac{2}{\mu(t + \gamma_1)}$.

Let $\Delta_t := \mathbb{E}[\mathcal{T}(\theta_{k,t}) - \mathcal{T}(\theta_k^*)]$. Then the inequality becomes:
\begin{equation}
\sum_{t=0}^{T-1} \frac{4}{\mu(t + \gamma_1)} \Delta_t \leq D + \sum_{t=0}^{T-1} \frac{4}{\mu^2(t + \gamma_1)^2} C,
\end{equation}
where $D = \|\theta_{k,0} - \theta_k^*\|^2$.

Next, apply Jensen’s inequality and convexity of $\mathcal{T}(\cdot)$:
\begin{equation}
\mathbb{E}[\mathcal{T}(\bar{\theta}_{k,T}) - \mathcal{T}(\theta_k^*)] \leq \frac{1}{T} \sum_{t=0}^{T-1} \Delta_t.
\end{equation}

We now bound the harmonic and squared harmonic series using integral approximation:
\begin{align}
    \sum_{t=0}^{T-1} \frac{1}{t + \gamma_1} &\leq \log\left(\frac{T}{\gamma_1} + 1\right), \\
    \sum_{t=0}^{T-1} \frac{1}{(t + \gamma_1)^2} &\leq \frac{1}{\gamma_1}.
\end{align}

Thus,
\begin{equation}
\sum_{t=0}^{T-1} \Delta_t \leq \frac{\mu D}{4} \log\left(\frac{T}{\gamma_1} + 1\right) + \frac{C}{\mu\gamma_1}.
\end{equation}

Dividing by $T$, we get:
\begin{equation}
\mathbb{E}[\mathcal{T}(\bar{\theta}_{k,T}) - \mathcal{T}(\theta_k^*)]
\leq \frac{\mu D}{4T} \log\left(\frac{T}{\gamma_1} + 1\right) + \frac{C}{\mu \gamma_1 T}.
\end{equation}

By absorbing constants into a single constant (e.g., involving $L_1$), we obtain:
\begin{equation}
\mathbb{E}[\mathcal{T}(\bar{\theta}_{k,T})] - \mathcal{T}(\theta_k^*)
\leq \frac{4L_1\gamma_1}{\mu T} \left(1 + \log\left(\frac{T}{\gamma_1} + 1\right)\right),
\end{equation}
which proves the desired bound for client models.
\end{proof}

\begin{proof}[Theorem~\ref{theo:1} (Convex \emph{PPGM} Bound)]
To derive the convergence for the \emph{PPGM} model, we start by summing the result of Lemma~\ref{lemma:2} over all $T$ rounds:

\begin{equation}
\label{eq:ppgm-1}
    \sum_{t=0}^{T-1} 2\eta_g \mathbb{E}[\mathcal{V}(\psi_{k,t}) - \mathcal{V}(\psi_k^*)] 
    \leq \mathbb{E}[\|\psi_{k,0} - \psi_k^*\|^2] + T\eta_g^2 \sigma_2^2.
\end{equation}

Using Jensen's inequality and convexity of $\mathcal{V}(\cdot)$:
\begin{equation}
\label{eq:ppgm-2}
    \mathbb{E}[\mathcal{V}(\bar{\psi}_{k,T}) - \mathcal{V}(\psi_k^*)] 
    \leq \frac{1}{T} \sum_{t=0}^{T-1} \mathbb{E}[\mathcal{V}(\psi_{k,t}) - \mathcal{V}(\psi_k^*)].
\end{equation}

Multiply both sides of Equation~\eqref{eq:ppgm-2} by $2T\eta_g$:
\begin{equation}
    2T\eta_g \mathbb{E}[\mathcal{V}(\bar{\psi}_{k,T}) - \mathcal{V}(\psi_k^*)] 
    \leq 2\eta_g \sum_{t=0}^{T-1} \mathbb{E}[\mathcal{V}(\psi_{k,t}) - \mathcal{V}(\psi_k^*)].
\end{equation}

Applying Equation~\eqref{eq:ppgm-1}, we get:
\begin{equation}
    2T\eta_g \mathbb{E}[\mathcal{V}(\bar{\psi}_{k,T}) - \mathcal{V}(\psi_k^*)] 
    \leq \mathbb{E}[\|\psi_{k,0} - \psi_k^*\|^2] + T\eta_g^2 \sigma_2^2.
\end{equation}

Divide both sides by $2T\eta_g$:
\begin{equation}
\label{eq:ppgm-intermediate}
    \mathbb{E}[\mathcal{V}(\bar{\psi}_{k,T}) - \mathcal{V}(\psi_k^*)] 
    \leq \frac{\mathbb{E}[\|\psi_{k,0} - \psi_k^*\|^2]}{2T\eta_g} + \frac{\eta_g \sigma_2^2}{2}.
\end{equation}

Now use the learning rate $\eta_g = \frac{2}{L_2(t+\gamma_2)}$ with $\gamma_2 = \max\left\{\frac{8L_2}{\mu}, E\right\}$. Since $\eta_g$ is non-increasing, we have:
\begin{equation}
    \eta_g \leq \frac{2}{L_2 \gamma_2}.
\end{equation}

Substituting this into Equation~\eqref{eq:ppgm-intermediate}:
\begin{equation}
    \mathbb{E}[\mathcal{V}(\bar{\psi}_{k,T}) - \mathcal{V}(\psi_k^*)] 
    \leq \frac{L_2 \gamma_2 \mathbb{E}[\|\psi_{k,0} - \psi_k^*\|^2]}{4T} + \frac{\sigma_2^2}{L_2 \gamma_2}.
\end{equation}

Now, applying the $L_2$-smoothness of $\mathcal{V}(\cdot)$:
\begin{equation}
    \mathbb{E}[\|\psi_{k,0} - \psi_k^*\|^2] \leq \frac{2}{L_2} (\mathcal{V}(\psi_{k,0}) - \mathcal{V}(\psi_k^*)).
\end{equation}

Thus, we get:
\begin{equation}
    \mathbb{E}[\mathcal{V}(\bar{\psi}_{k,T}) - \mathcal{V}(\psi_k^*)] 
    \leq \frac{\gamma_2 (\mathcal{V}(\psi_{k,0}) - \mathcal{V}(\psi_k^*))}{2T} + \frac{\sigma_2^2}{L_2 \gamma_2}.
\end{equation}

Finally, this bound can be simplified using the inequality $\frac{\gamma_2}{2T} \leq \frac{4\gamma_2}{T}(1 + \log(\frac{T}{\gamma_2} + 1))$ for $T \geq 1$, and absorbing constants:
\begin{equation}
\begin{split}
    \mathbb{E}[\mathcal{V}(\bar{\psi}_{k,T})] - \mathcal{V}(\psi_k^*) 
    \leq \frac{4L_2\gamma_2}{T} \left(1 + \log\left(\frac{T}{\gamma_2} + 1\right)\right),
\end{split}
\end{equation}
\end{proof}


\begin{theorem}[Non-Convex Convergence Bound for Client Model]
\label{theo:2}
Let Assumptions 1--4 hold. Suppose a fixed client \( k \) performs local SGD with step size \( \eta_l \) for \( T \) rounds, each consisting of \( E \) local updates. Then, the iterates \( \{\theta_{k,t}\}_{t=0}^{T-1} \) satisfy:
\begin{equation}
    \sum_{t=0}^{T-1} \mathbb{E}[\|\nabla \mathcal{T}(\theta_{k,t})\|^2]
    \leq \frac{2}{\eta_l} \left( \mathcal{T}(\theta_{k,0}) - \mathcal{T}(\theta_k^*) \right) + L_1 T \eta_l \sigma_1^2.
\end{equation}
% where \( \gamma_1 = \max\left\{ \frac{8L_1}{\mu}, E \right\} \).
\begin{equation}
	\sum_{t=0}^{T-1} \mathbb{E}[\|\nabla \mathcal{V}(\psi_{k,t})\|^2]
	\leq \frac{2}{\eta_g} \left( \mathcal{V}(\psi_{k,0}) - \mathcal{V}(\psi_k^*) \right) + L_2 T \eta_g \sigma_2^2.
\end{equation}
\end{theorem}

\begin{proof}[Theorem~\ref{theo:2} (Non-Convex Client Model Bound)]
\label{proof:theo2-1}
The local SGD update rule for client \( k \) at round \( t \) is:
\begin{equation}
    \theta_{k,t+1} = \theta_{k,t} - \eta_l g_{k,t},
\end{equation}
where \( g_{k,t} \) is the stochastic gradient estimator. Under the \( L_1 \)-smoothness of \( \mathcal{T}(\theta_k) \), we have:
\begin{equation}
    \mathcal{T}(\theta_{k,t+1}) \leq \mathcal{T}(\theta_{k,t}) - \eta_l \langle \nabla \mathcal{T}(\theta_{k,t}), g_{k,t} \rangle + \frac{L_1 \eta_l^2}{2} \|g_{k,t}\|^2.
\end{equation}
Taking expectation conditioned on \( \theta_{k,t} \) and using \( \mathbb{E}[g_{k,t}] = \nabla \mathcal{T}(\theta_{k,t}) \):
\begin{equation}
\begin{split}
    \mathbb{E}[\mathcal{T}(\theta_{k,t+1})]
    &\leq \mathcal{T}(\theta_{k,t}) - \eta_l \|\nabla \mathcal{T}(\theta_{k,t})\|^2 + \frac{L_1 \eta_l^2}{2} \mathbb{E}[\|g_{k,t}\|^2].
\end{split}
\end{equation}
Using the variance bound \( \mathbb{E}[\|g_{k,t}\|^2] \leq \|\nabla \mathcal{T}(\theta_{k,t})\|^2 + \sigma_1^2 \), we get:
\begin{equation}
\begin{split}
    \mathbb{E}[\mathcal{T}(\theta_{k,t+1})]
    &\leq \mathcal{T}(\theta_{k,t}) - \left(\eta_l - \frac{L_1 \eta_l^2}{2}\right) \|\nabla \mathcal{T}(\theta_{k,t})\|^2\\ 
    & \qquad \qquad \qquad+ \frac{L_1 \eta_l^2}{2} \sigma_1^2.
\end{split}
\end{equation}
Summing over \( t = 0 \) to \( T-1 \), we obtain:
\begin{equation}
\begin{split}
    \sum_{t=0}^{T-1} \left( \eta_l - \frac{L_1 \eta_l^2}{2} \right) \mathbb{E}[\|\nabla \mathcal{T}(\theta_{k,t})\|^2]
    &\leq \mathcal{T}(\theta_{k,0}) - \mathcal{T}(\theta_k^*) \\
    & \quad + \frac{L_1 T \eta_l^2}{2} \sigma_1^2.
\end{split}
\end{equation}

Assuming \( \eta_l \leq \frac{1}{L_1} \), we simplify the left-hand side:
\begin{equation}
    \sum_{t=0}^{T-1} \mathbb{E}[\|\nabla \mathcal{T}(\theta_{k,t})\|^2]
    \leq \frac{2}{\eta_l} \left( \mathcal{T}(\theta_{k,0}) - \mathcal{T}(\theta_k^*) \right) + L_1 T \eta_l \sigma_1^2.
\end{equation}

\end{proof}

\begin{proof}[Theorem~\ref{theo:2} (Non-Convex PPGM Bound)]
	Identical in structure to Proof~\ref{proof:theo2-1}, replacing $T,\theta,\eta_l,L_1,\sigma_1^2$ 
	with $V,\psi,\eta_g,L_2,\sigma_2^2$.
\end{proof}


\subsection{Privacy Analysis of \textsf{ProtoGen}}

Let each client maintain a feature extractor $\mathcal{H}(\phi_k,x)$ and a \emph{PPGM} model with encoder $\mathcal{E}(\rho_k,z|y)$ and decoder $\mathcal{J}(\varrho_k, e|y)$. The privacy mechanism relies on input perturbation, followed by feature extraction and training a generator only on the perturbed representations.

\begin{assumption}[Input Perturbation by Gaussian Noise]
\label{asm:6}
Each input $x_i$ is perturbed using additive Gaussian noise:
\[
\tilde{x}_i = x_i + \xi_i, \quad \text{where } \xi_i \sim \mathcal{N}(0, \sigma^2 I_d).
\]
\end{assumption}

\begin{assumption}[Lipschitz Continuity of the Feature Extractor]
\label{asm:7}
The feature extractor $\mathcal{H}$ is $L$-Lipschitz continuous, i.e., for all $x, x' \in \mathbb{R}^d$,
\[
\|\mathcal{H}(\phi_k,x) - \mathcal{H}(\phi_k,x')\|_2 \leq L \cdot \|x - x'\|_2.
\]
\end{assumption}

\begin{definition}[Differential Privacy~\cite{dwork2014algorithmic}]
A randomized algorithm $\mathcal{M}$ satisfies $(\varepsilon, \delta)$-differential privacy if for any two adjacent datasets $\mathcal{D}, \mathcal{D}'$ differing in at most one record, and for any measurable subset $S$ of outputs,
\[
\Pr[\mathcal{M}(\mathcal{D}) \in S] \leq e^\varepsilon \Pr[\mathcal{M}(\mathcal{D}') \in S] + \delta.
\]
\end{definition}

\begin{definition}[Global Sensitivity]
The $L_2$-sensitivity of a function $\mathcal{H}(\phi_k,\mathcal{X}_k) \rightarrow \mathcal{Z}_k$ is
\[
\Delta_h = \sup_{\mathcal{D}, \mathcal{D}' \text{ differ in one sample}} \|\mathcal{H}(\phi_k,\mathcal{D}) - \mathcal{H}(\phi_k,\mathcal{D}')\|_2.
\]
\end{definition}

\begin{theorem}[Gaussian Mechanism~\cite{dwork2014algorithmic}]
Let $\mathcal{H}(\phi_k,\mathcal{X}_k) \rightarrow \mathcal{Z}_k$ be a function with $L_2$ sensitivity $\Delta_h$. Then the mechanism
\[
\mathcal{M}(\mathcal{D}) = \mathcal{H}(\phi_k,\mathcal{D}) + \mathcal{N}(0, \sigma^2 I_d)
\]
is $(\varepsilon, \delta)$-differentially private for any $\delta > 0$, where
\[
\varepsilon = \frac{\Delta_f^2}{2 \sigma^2} + \log(1/\delta).
\]
\end{theorem}

\begin{lemma}[Differential Privacy of Noisy Feature Extraction]
\label{lemma:3}
Let $\mathcal{M}(x) = \mathcal{H}(\phi_k, x + \xi)$ where $\xi \sim \mathcal{N}(0, \sigma^2 I_d)$, and let the input sensitivity $\Delta_x := \sup_{x,x'} \|x - x'\|_2$ be bounded. Then $\mathcal{M}$ satisfies $(\varepsilon, \delta)$-differential privacy with
\[
\varepsilon = \frac{L^2 \Delta_x^2}{2 \sigma^2} + \log(1/\delta).
\]
\end{lemma}

\begin{proof}
We consider two neighboring datasets that differ in one input point $x_i$ versus $x_i'$. We analyze the output of the mechanism $\mathcal{M}(x) = \mathcal{H}(\phi_k, x + \xi)$. Since $\mathcal{H}$ is $L$-Lipschitz and $\xi$ is sampled independently from $\mathcal{N}(0, \sigma^2 I_d)$:

\begin{equation}
\|\mathcal{M}(x) - \mathcal{M}(x')\|_2 
= \|\mathcal{H}(\phi_k, x + \xi) - \mathcal{H}(\phi_k, x' + \xi)\|_2
\end{equation}

Using assumption \ref{asm:7}, we have following inequality:
\begin{equation}
\|\mathcal{M}(x) - \mathcal{M}(x')\|_2  \leq L \cdot \|x - x'\|_2.
\end{equation}

Hence, the sensitivity of the mapping $x \mapsto \mathcal{H}(\phi_k, x + \xi)$ is at most $L \cdot \Delta_x$. By applying the Gaussian mechanism, this perturbation ensures $(\varepsilon, \delta)$-DP for
\[
\varepsilon = \frac{(L \Delta_x)^2}{2 \sigma^2} + \log(1/\delta).
\]
\end{proof}

\begin{lemma}[Post-Processing Invariance~\cite{dwork2014algorithmic}]
\label{lemma:4}
Let $\mathcal{M} : \mathcal{D} \rightarrow \mathcal{Y}$ be an $(\varepsilon, \delta)$-DP mechanism. Let $f : \mathcal{Y} \rightarrow \mathcal{Z}$ be any (possibly randomized) function. Then $f \circ \mathcal{M}$ also satisfies $(\varepsilon, \delta)$-DP.
\end{lemma}

\begin{theorem}[Privacy Guarantee of \emph{PPGM} Training and Sharing]
\label{theo:4}
Let each feature $z_i$ used for \emph{PPGM} model training be generated via $z_i = \mathcal{H}( \phi_k, x_i + \xi_i)$, where $\xi_i \sim \mathcal{N}(0, \sigma^2 I_d)$ and $\mathcal{H}$ is $L$-Lipschitz. Then:

\begin{enumerate}
    \item The feature memory $\tilde{\mathcal{D}}_k = \{(\tilde{z}_i, y_i)\}_{i=1}^m$ satisfies $(\varepsilon, \delta)$-differential privacy with respect to the original data $\sD_k$.
    \item The decoder $\mathcal{J}_{\varrho_k}$ trained on $\tilde{\mathcal{D}}_k$ also satisfies $(\varepsilon, \delta)$-DP due to post-processing.
\end{enumerate}
\end{theorem}

\begin{proof}[Theorem \ref{theo:4}]

\textbf{Step 1: Privacy of Feature Extraction.}  
From Lemma \ref{lemma:3}, each feature $\tilde{z}_i = \mathcal{H}(\phi_k, x_i + \xi_i)$ is $(\varepsilon, \delta)$-DP due to the Gaussian mechanism applied to the input $x_i$, and the Lipschitz property of $\mathcal{H}$.

\textbf{Step 2: Privacy of \emph{PPGM} model Training.}  
The \emph{PPGM} is trained on the dataset $\tilde{\mathcal{D}}_k$ composed of $(\tilde{z}_i, y_i)$. Since the entire training procedure (gradient descent) operates on already private representations, and no raw $x_i$ or unperturbed $z_i$ are used, by Lemma \ref{lemma:4} (post-processing), the resulting trained model — including the decoder $\mathcal{J}(\varrho_k,e)$ — is also $(\varepsilon, \delta)$-DP.

\textbf{Step 3: Sharing the Decoder.}  
Only the decoder parameters $\varrho_k$ are transmitted to the server. These parameters cannot be reversed for any particular $\tilde{z}_i$ due to aggregated training and noise injection. Because $\mathcal{J}(\varrho_k,e)$ is computed via post-processing of DP data, it retains the same $(\varepsilon, \delta)$ privacy guarantees.

Thus, the privacy of any individual data sample in the client dataset $\sD_k$ is preserved throughout the entire learning procedure.
\end{proof}




\bibliographystyle{ieeetr}
\bibliography{reference}



\end{document}
